\documentclass{plantilla}

% --- Metadatos del Informe ---
\titulo{Plantilla para informes en fomato IEEE}
\autores{Pedro Ernst, Valentino Rao}
\materia{Nombre de la materia}
\comision{4R2}

\begin{document}

\maketitle

\begin{abstract}
En este informe se describe es un ejemplo de como utilizar una plantilla en formato cls para la redaccion de informes en formto IEEE.
\end{abstract}

\begin{IEEEkeywords}
Formato IEEE, plantilla.
\end{IEEEkeywords}

\section{Introducción}
\IEEEPARSTART{E}{l} En el siguiente informe analizaremos...

\section{Desarrollo Experimental}
Para el cálculo de las resistencias de polarización, utilizamos las unidades que definimos en la plantilla:
$R_1 = 10\kohm$ y $C_1 = 100\uF$.

\begin{equation}
    V_{out} = A_v \cdot V_{in}
\end{equation}

\section{Resultados}

\begin{figure}[H]
\centering
\includegraphics[width=6.5cm]{images/circuito.png}
\caption{Imagen del circuito implementado.}
\label{fig_circuito}
\end{figure}

\section{Conclusión}
Se logró verificar que...

% Referencias bibliográficas
\begin{thebibliography}{1}
\bibitem{IEEEhowto:kopka}
H. Kopka and P. W. Daly, \emph{A Guide to \LaTeX}, 3rd ed.\hskip 1em plus
  0.5em minus 0.4em\relax Harlow, England: Addison-Wesley, 1999.
Bubylibro Medios de enlace, La vida hecha libro. 
\bibitem{IEEEhowto:kopka}
Bubylibro Medios de enlace, La vida hecha libro. 
\end{thebibliography}

\end{document}